% Short running form; the full title is too wide for the head's right box.
\runninghead{The Peace of a City}
\properlane{peace-of-the-house}{The Peace of a City Still Being Built: From ``Da pacem'' to the Courts of the Lord}
\label{sec:peace-of-the-house}

Heard from its chants, the Mass is one ascent from a request to an entry. The
Introit asks for peace and names a destination, ``We shall go into the house of
the Lord''. The Gradual sings the same words again and prays for the peace of
the city whose walls and towers it names. The Alleluia sees the nations and
their kings coming to fear the Lord's name as Sion is built. The Communion
tells the worshippers to bring up sacrifices and come into his courts. The
peace for which the Introit asks is therefore not the mere absence of trouble.
It is the completion of a building: the house of the Lord to which the people
go up, the city whose strength that peace is, and the courts into which, at
the end of the Mass, they are told to enter. St Augustine, expounding all three
psalms, St Hilary of Poitiers on the Gradual's, and Cassiodorus beside them
read the house and the city as the Church being built of living stones and
completed in heaven, and they say what it is built of. St John Chrysostom,
expounding the Epistle, says what threatens it.

\subsection{Going up to the house: Ps 121:1 in Augustine, Hilary and Cassiodorus}

St Augustine begins his exposition of Ps~121 from its title, \latin{Canticum
graduum}, a Song of degrees, ``for these degrees are not of descent, but of
ascent''. What the singer ascends to is not the sky but a city: ``there is in
heaven the eternal Jerusalem, where are our fellow-citizens, the Angels: we are
wanderers on earth from these our fellow-citizens. We sigh in our pilgrimage;
we shall rejoice in the city.'' Then he gives the first verse its point: ``we
find companions in this pilgrimage, who have already seen this city herself;
who summon us to run towards her. At these he also rejoiceth, who saith, `I
rejoiced in them who said unto me, We will go into the house of the Lord'\,''
(\work{Enarrationes in Psalmos} 121, 2). Augustine's psalter did not read the
verse as the Missal sings it. Where the Missal has \latin{Laetatus sum in his,
quae dicta sunt mihi}, joy in the things that were said, his text read
\latin{Jucundatus sum in his qui dixerunt mihi}, joy in those who said, and he
builds on the persons: \latin{Jucundatus sum in Prophetis, jucundatus sum in
Apostolis. Omnes enim isti dixerunt nobis, In domum Domini ibimus}. The
reading therefore belongs to Augustine's words and not to the Missal's; what it
brings to the Mass is the plural of the verse's second half, \latin{ibimus},
which the Missal and Augustine share. No one goes up to this house alone.

St Hilary of Poitiers reaches the same city by another road. The psalm, he
says, holds nothing obscure for one who is held by the hope of it: \latin{Qui
spe caelestis desiderii detinetur, nihil obscuritatis in psalmo habebit}. A man
who remembers that he is coheir, \latin{futurumque incolam ciuitatis caelestis
uiuis lapidibus extructae}, a future dweller in the heavenly city built of
living stones, must cry with the prophet \latin{Laetatus sum in his, quae dicta
sunt mihi} (\work{Tractatus super Psalmos}, in Ps.~121, 1). The house is the
Sion that the Lord chose for his rest, and Hilary refuses the earthly city in
so many words: \latin{non utique \dots\ hanc terrenam et caducam \dots\ sed illam
liberam et caelestem Hierusalem, quia eadem et Sion est} (121, 2). He names its
builders: \latin{a Moyse in lege, a prophetis in passionibus suis, a domino in
corpore, ab apostolis in martyriis, a sancto spiritu in uirtutibus. hi
aedificatores, haec aedificia, haec ciuitas} (121, 3). The psalm's ``built as a
city'' means what it says: \latin{non ciuitas, sed ut ciuitas}, not a city
but as a city, because the building of the earthly city, the raising of the
Temple and the setting up of the tabernacle foreshadowed the eternal and
heavenly one, which is being built still and so is given no date (121, 4).

What Hilary adds to Augustine is the reason for the plural. We go up together
because the Church is one body: \latin{quia unum ecclesiae corpus est, non
quadam corporum confusione permixtum \dots\ sed per fidei unitatem, per
caritatis societatem, per operum uoluntatisque concordiam, per sacramenti unum
in omnibus donum, unum omnes sumus} (121, 5). The unity he describes is not a
heap of bodies thrown together but the oneness of faith, charity, common work
and the one gift of the sacrament in all. Among his proofs he sets Paul's plea
to the Corinthians as his text gives it, that all be of one mind, exercising
the same charity (121, 5). Zingerle's edition refers the plea to 1 Cor 1:10,
two verses after this Sunday's Epistle ends; Paul's own verse there also asks
``that there be no schisms among you''.

Cassiodorus, reading the same verse, asks who said the words at which the
psalmist rejoices, and answers \latin{Scilicet Spiritus sanctus, qui cordi eius
tacita voce loquebatur}: the Holy Spirit, speaking to his heart with a silent
voice. The house is \latin{domus desiderabilis, domus de vivis lapidibus
fabricata}, a house to be longed for, made of living stones, and the joy is
worthy of it: \latin{O digna exsultatio illuc ire contendere, unde nunquam
aliquis velit exire!} (\work{Expositio psalmorum}, in Ps.~121). The three
agree on what matters to this reading: the house to which the Introit and the
Gradual go up is a city of persons, built of those who go up to it.

\subsection{Peace in thy strength: what the city is built of}

The Gradual's versicle prays \latin{Fiat pax in virtute tua}, ``Let peace be in
thy strength''. Augustine asks what a city's strength is and answers with
charity: ``O Jerusalem, O city, who art being built as a city, whose partaking
is in `The Same:' `Peace be in thy strength:' peace be in thy love; for thy
strength is thy love.'' He proves it from the Song of Songs, ``Love is strong
as death'', and draws the consequence for the one who is built into the city:
``since this love slayeth what we have been, that we may be what we were not;
love createth a sort of death in us'' (\work{Enarrationes in Psalmos} 121, 12).
The peace the Introit asks for is thus received as charity, and charity is
costly: it puts to death the self that could not be at peace. Cassiodorus
reaches the same identification in almost the same words: \latin{Virtus quippe
ipsius pax sine dubitatione sanctorum est, quae vocatur et charitas}; by it the
saints \latin{fiunt unum}, and by it they deserve to be the Creator's temple
(\work{Expositio psalmorum}, in Ps.~121).

Hilary reads the verse otherwise. For him
peace and strength are joined but not identical: \latin{pax atque uirtus domus
huius conexa sunt sibi}, and nothing can be strong except out of peace (in
Ps.~121, 13). And again: \latin{pax enim ecclesiae et unitas uirtutem nobis
firmitatemque praestat}, the Church's peace and unity give us strength and
firmness (121, 14). For Augustine and Cassiodorus the city's strength is its charity; for Hilary it is
the firmness that its peace confers. On the towers of the same verse the three
part further. Cassiodorus reads them as the martyrs, who defend the city of God
by setting their bodies against its enemies in faithful confession. Hilary
reads them as the faithful and holy men who give the rest of the city the
firmness of towers (121, 14). Augustine, in the Latin that the English
translation of his sermon abridges, reads them as the city's \latin{excelsi}:
\latin{Pauci enim sedebunt in judicio; sed multi ad dexteram positi facient
populum civitatis illius}, the few who will sit in judgment, to whom the many
at the right hand belong. The three witnesses give the same clause three
allegories.

\subsection{The second wall: the nations of the Alleluia}

The Alleluia's verse, ``The Gentiles shall fear thy name, O Lord'', is for
Augustine the moment when a second wall arrives. He reads it in the order of
the psalm: God has had mercy on Sion, his servants have loved her stones, and
therefore the nations come. ``Now that Thou hast pitied Sion, now that Thy
servants have taken pleasure in her stones, by acknowledging the foundation of
the Apostles and Prophets \dots\ hence preaching hath increased among the
heathen: let the heathen fear Thy Name, let another wall approach also from the
heathen, let the Corner Stone be recognised, let the two who come from
different regions, but who no longer differ in belief, meet in close union''
(\work{Enarrationes in Psalmos} 101, s.~1, 16). The next verse, which the
Alleluia does not sing, gives the work its present tense: ``This work is going
on now. O ye living stones, run to the work of building, not to ruin'' (101,
s.~1, 17). Cassiodorus names the same building from the same verse:
\latin{aedificata est Sion, hoc est mater Ecclesia, de vivis lapidibus
fabricata} (\work{Expositio psalmorum}, in Ps.~101).

At this psalm the Greek exegesis joins the Latin. Theodoret, who reads Ps~121
of the earthly city alone, reads Ps~101:16 first of the neighbours who would
marvel at Israel's return from Babylon, and then says that it was fulfilled
properly and truly only after the Incarnation of our God and Saviour, when the
nations and their kings were converted (\work{Interpretatio in Psalmos}, in
Ps.~101; PG 80, cols.~1679--1680). The Greek expositions printed in Migne under
the name of St Athanasius join the calling of the nations to the stones of the
former people and read the Sion of the next verse as the Church (PG 27,
col.~429).

\subsection{Into the courts: the Communion}

The Communion's psalm carries the building in its title, ``when the house was
built after the captivity'', and Augustine reads the whole psalm under it. The
house of the title, he tells his hearers, is no house of hewn stone and raised
columns: ``It is no such house that is in building; for behold where it is
built, not in one spot, not in any particular region'' (\work{Enarrationes in
Psalmos} 95, 1). At the Communion's own verse he describes the entry: ``Behold
the house increasing: behold the edifice pervade the whole world. Rejoice,
because ye have entered into the courts; rejoice, because ye are being built
into the temple of God. For those who enter are themselves built up, they
themselves are the house of God'' (95, 9). And of the second clause he says
plainly which court is meant: ``O worship the Lord in His holy court: in the
Catholic Church; this is His holy court'' (95, 10).

Cassiodorus reads the title on two levels at once. According to the letter it
signifies the rebuilding of the Temple at Jerusalem by Zorobabel after the
captivity; according to the spirit, \latin{Domus enim ista, id est universalis
Ecclesia, in qua Christus inhabitat, vivis lapidibus semper exstruitur: quia
quotidie de confitentibus aedificationis sumit augmentum; nec aedificari
desinit donec usque ad finem saeculi praedestinatorum numerus impleatur}. The
Church in which Christ dwells is always being built of living stones, growing
every day by those who confess him, and it will not cease to be built until the
number of the predestined is filled at the end of the age (\work{Expositio
psalmorum}, in Ps.~95). With that sentence Cassiodorus uses one image, the
house of living stones, at the Gradual's psalm, at the Alleluia's and at the
Communion's.

Set in the order of the Mass, the three psalms describe one movement. The
Introit and the Gradual go up to the house; the Alleluia sees the second wall
come from the nations; the Communion enters the courts, and those who enter are
themselves built into the house. The request of the first words, \latin{Da
pacem}, is answered in the chants by the completion of a building whose
strength is charity.

\subsection{What threatens the city: the Epistle, and Christ's own city}

The Epistle shows a church that is not at peace. The Corinthians are rich ``in
all utterance and in all knowledge'', and, as the verses after the lesson
reveal, they are dividing into parties. St John Chrysostom reads Paul's
thanksgiving with that division in view. Under the colour of praise, he says,
Paul ``touches them sharply'', and his remedy is a name repeated until it
binds: ``consider how he always fastens them as with nails to the Name of
Christ'', so that ``by incessant application of that glorious Name he may
foment their inflammation, and purge out the corruption of the disease''
(\work{Homilies on First Corinthians} 2). A city is not held together by its
gifts. It is held together by the one Name its members bear in common.

Augustine and Chrysostom agree that the city's unity is not its citizens'
achievement: for Augustine the two walls meet in the Corner Stone, for
Chrysostom the parties are nailed to one Name. They differ in what they are
looking at. Augustine watches a building rise; Chrysostom attends to one that
is coming apart. Chrysostom's Corinth is what Augustine's Jerusalem is not
yet, and the Mass sets the lesson of a divided
church between two chants of the city at peace.

The Gospel's ``his own city'' enters this reading through St Hilary's
commentary on Matthew. Hilary reads the Gerasene episode and the healing of the
paralytic as one figure. The town that came out to meet Christ and begged him
to leave its borders \latin{Judaici populi habet speciem}, bears the likeness
of the Jewish people, for the law does not receive the Gospels: \latin{a qua
repudiatus, in civitatem suam revertitur \dots\ Deo civitas fidelium plebs est.
In hanc igitur navi, id est Ecclesia, vectus introiit}. Rejected by that town,
Christ returns to his own city, and to God the city is the people of the
faithful; he entered it carried in the boat, which is the Church
(\work{Commentarius in Matthaeum} VIII.4; PL 9, cols.~959--960). Hilary names
no town, so the allegory stands beside any literal identification of the place.
Aquinas, in his lectures on Matthew, gives the same kind of reading beside his
harmony of the Gospels: \latin{Et venit in civitatem suam,
scilicet in civitatem gentium, quae sibi datae sunt}, with the promise of
Ps~2:8, ``I will give thee the Gentiles for thy inheritance'' (\work{Super
Matthaeum} c.~9). Read with Hilary and Aquinas, the boat that crosses to the
city in the Gospel and the house that the chants go up to are the same Church
seen from two sides.

\subsection{The prayers and the altar}

The orations speak in the plural of a people. The Introit's antiphon ends with
\latin{plebis tuae Israel}; the Collect asks that God direct \latin{corda
nostra}, and gives the reason that no one who goes up can supply: \latin{tibi
sine te placere non possumus}. The ascent is not self-propelled. The Offertory
ends where the city's life begins, \latin{in conspectu filiorum Israel}: Moses
offers, but the people are present at their own sacrifice, and those closing
words are the one clause of the antiphon that stands verbatim in the chapter it
cites, at 24:17. The Secret names the unity toward which the city is built, and
names it first as God's: \latin{unius summae divinitatis participes}, sharers
in the one supreme Godhead. The Postcommunion is the thanksgiving of those who
have entered the courts, and it still asks, \latin{ut dignos nos eius
participatione perficias}, that they be made worthy of what they share: the
courts have been entered, and the city is still ahead.

\subsection{The earthly city rebuilt, the Donatists and the four chants}

The Gradual's psalm, on which this reading leans most, has another reading
among its Greek interpreters, and it is not an allegory. St John Chrysostom
reads Ps~121 of the exiles returning from Babylon. The captivity made them
better; they longed for the house of prayer and the city; and at verse 3 the
city built as a city is the city they found in ruins, its towers thrown down
and its walls cast to the ground (\work{Expositio in Psalmum CXXI}, PG 55,
cols.~347--348). At the Gradual's versicle he explains ``strength'' as the
city's substance, its inhabitants and its plenty, and says the psalmist prays
peace for her because war had destroyed her (col.~350). Theodoret reads the
psalm the same way: the returning exiles rejoice not as men about to recover
their houses but as men about to see the house of God, and the prayer of verse
7 is for the city's walls, palaces and houses (PG 80, cols.~1879--1882).
Neither reads the Church or heaven at this psalm. Chrysostom's own turn to his
hearers is moral, and it serves this reading's literal sense well: call men to
the hippodrome and many run together, but call them to the house of prayer and
few are not sluggish, so that Christians show themselves slower to go up than
the returning Jews were.

The return from exile is therefore the literal sense of the Gradual, and the
heavenly city of Augustine, Hilary and Cassiodorus its allegorical sense. St Robert Bellarmine holds both in that order: the psalm is
``the language of God's people, expressive of their joy on hearing the welcome
news of their return to their country'', but ``Christ, however, was the bearer
of a far and away more happy message'', and the psalm ``treats of the celestial,
and not the earthly Jerusalem'' (\work{A Commentary on the Book of Psalms}, on
Ps~121). So does Schuster, commenting on this Mass: the psalmist rejoices
``after the afflictions of the Babylonian exile'', and then ``All this is, of
course, to be glorified by a spiritual interpretation. The peace which is here
described is the atmosphere of the heavenly Jerusalem.'' The medieval
commentators on the Mass carry the same division into the chant. Sicard, and
Honorius in the first of his two readings, hear its chants of the people's
return from captivity; in his second, Honorius hears St Gregory the Great pray
in the Gradual's verse \latin{ut pax in turribus Ecclesiae fiat}, that there be
peace in the Church's towers.

A second difficulty lies in Augustine's sermons themselves. Their sharpest
paragraphs belong to his controversy with the Donatists and turn the psalms
against those who ``say, Peace be with you: and have not the
peace which they preach to the people'' (121, 13). Three of their
propositions do not depend on that setting: the city's peace is its charity,
the courts into which the faithful enter are one, and those who enter are
built into the house.

A third objection is that the chants are simply four texts, each chosen for a
word, \latin{pacem}, \latin{domum}, \latin{gentes}, \latin{hostias}, and that
they describe no movement at all. One fact about the formulary tells
against that account: Ps~121:1 is sung twice, as the Introit's verse and as the
Gradual's respond, so that the destination named at the entrance is named again
after the Epistle.

\subsection{The four senses of this reading}

\begin{fourSenses}
\item[Literal.] A people waiting on the Lord prays for peace and for its
prophets to be found faithful. Pilgrims rejoice at the call to go up to the
house of the Lord in Jerusalem, and pray for peace within the walls and
abundance in the towers of a city that the exiles found in ruins
(Chrysostom, Theodoret). A psalm written for the house built after the
captivity summons the nations to bring their sacrifices into its courts, and
another foretells that the nations and their kings will fear the Lord's name
when he has built up Sion.

\item[Allegorical.] The city is the Church, built of living stones
(Augustine, Hilary, Cassiodorus). Its builders are the law, the prophets,
Christ in his body, the apostles and the Spirit (Hilary); its two walls, from
Israel and from the nations, meet in the Corner Stone (Augustine); its strength
is charity (Augustine, Cassiodorus), or the firmness its peace and unity give
(Hilary); and its holy court is the Catholic Church (Augustine). Christ enters
it carried in the boat of the Church, and the city he returns to, rejected by
the town that figures the law, is the people of the faithful (Hilary, on
Matthew).

\item[Moral.] Peace is received as charity, and charity ``slayeth what we have
been, that we may be what we were not'' (Augustine). No one goes up alone: the
plural \latin{ibimus} is the unity of faith, charity and the one sacrament
(Hilary), and the gifts that divide a church are held together only by the Name
of Christ (Chrysostom). Christians who are slower to the house of prayer than
the exiles were to Jerusalem are to go up (Chrysostom), and the living stones
are to ``run to the work of building, not to ruin'' (Augustine).

\item[Anagogical.] The house the Introit says we will go into is entered in this
Mass in sign; the psalm's \latin{ibimus} stays future. The Communion's court is
the city's forecourt, and the city itself is the eternal Jerusalem where the
Angels are fellow-citizens (Augustine), \latin{unde nunquam aliquis velit
exire} (Cassiodorus). The building goes on until the number of the elect is
filled (Cassiodorus), and ``when the building is finished, and the house
dedicated'' the time for running is over (Augustine, on Ps~101).
\end{fourSenses}
